\subsubsection{Quantum Error Correction Codes (QECC)}\label{sec:qecc}

After understanding that physical qubits suffer from decoherence and gate errors, we need a \textbf{mechanism to preserve quantum information during a computation}. \hl{Quantum Error Correction Codes (QECCs)} achieve this by introducing \textbf{redundancy}: instead of storing information in a single physical qubit, the information is distributed across many qubits.

\highspace
We previously discussed this concept on \autopageref{sec:logical-vs-physical-qubits}, where we compared logical and physical qubits. Here, however, we present the concept with more detail because we now have a better understanding of possible errors and their impact on quantum computations.
\begin{itemize}
    \item A \definition{Physical Qubit} is a \textbf{real qubit implemented by hardware}, such as superconducting qubits, trapped-ion qubits or photonic qubits. Unfortunately, physical qubits are directly affected by decoherence (\autopageref{def:coherence}), thermal noise (\autopageref{sec:quantum-computation-errors}), and gate imperfections (infidelity, \autopageref{sec:gate-infidelity}). Therefore, a \textbf{physical qubit is unreliable if used alone} for quantum computation, as it is prone to errors that can lead to incorrect results.


    \item To obtain reliable computation, quantum information is not stored in a single physical qubit. Instead, one creates a \definition{Logical Qubit} $\ket{q_L}$ which is \textbf{encoded into many physical qubits}. For instance, a logical qubit can be encoded into 5, 7, or even hundreds of physical qubits, depending on the specific error correction code used.
    
    The user of the quantum computer thinks in terms of logical qubits, while the hardware manipulates the underlying physical qubits.
\end{itemize}

\begin{flushleft}
    \textcolor{Red2}{\faIcon{exclamation-triangle} \textbf{How can errors be detected without collapsing the state?}}
\end{flushleft}
Suppose our logical qubit is encoded into some physical qubits:
\begin{equation*}
    q_0, \, q_1, \, q_2, \, q_3, \, q_4
\end{equation*}
An error occurs on one of the physical qubits, say $q_2$. However, since the logical qubit is encoded across all five physical qubits, enough information remains in the qubits to recover the original state. The question is, \emph{\textbf{how do we know that an error occurred and which qubit was affected?}}

\highspace
\textcolor{Red2}{\faIcon{times-circle}} The first naïve solution could be to measure all the physical qubits to check for errors. However, this approach is flawed because \textbf{measuring a qubit collapses its quantum state}, which would destroy the quantum information we are trying to protect.

\highspace
At this point, we can \textbf{exploit the redundancy introduced by the logical encoding}. Consider a simple example where:
\begin{equation*}
    \ket{0_L} = \ket{000}
    \qquad
    \ket{1_L} = \ket{111}
\end{equation*}
Where $L$ stands for ``logical''. In this encoding, the logical qubit $\ket{0_L}$ is represented by three physical qubits all in the state $\ket{0}$, and the logical qubit $\ket{1_L}$ is represented by three physical qubits all in the state $\ket{1}$.
A generic logical qubit is therefore:
\begin{equation*}
    \ket{\psi_L}
    =
    \alpha\ket{000}
    +
    \beta\ket{111}
\end{equation*}
Notice that in a \textbf{valid encoded state all physical qubits agree with each other}: they are either all $0$ or all $1$.

\highspace
Now, suppose that a \textbf{bit-flip error} affects the second physical qubit. The error acts on the entire quantum state:
\begin{equation*}
    \alpha\ket{000}
    +
    \beta\ket{111}
    \quad\longrightarrow\quad
    \alpha\ket{010}
    +
    \beta\ket{101}
\end{equation*}
In this case, the second qubit has flipped from $0$ to $1$ in the first term and from $1$ to $0$ in the second term. The resulting state is \textbf{no longer a valid encoded state} because the physical qubits do not all agree with each other.

\highspace
\textcolor{Green3}{\faIcon{check-circle} \textbf{Solution: Syndrome.}} Rather than asking ``\emph{what is the quantum state?}'' which would destroy the computation, we only ask: ``\emph{are neighbouring qubits equal or different?}''. For instance, we can check if $q_0$ and $q_1$ are the same, and if $q_1$ and $q_2$ are the same. This way, we can detect that $q_2$ is different from the others without measuring the actual state of any qubit. In other words, the \textbf{pattern of answers identifies the location of the error without revealing the coefficients} $\alpha$ and $\beta$ of the logical qubit. The collection of these answers is called the \textbf{syndrome}.

\highspace
\begin{definitionbox}[: Syndrome]
    A \definition{Syndrome} is a classical piece of information that indicates the \textbf{presence and type of error without revealing the logical quantum state}. A \emph{syndrome} is obtained by performing a \definition{Syndrome Measurement}, which simply collects information about the error, a sort of ``\emph{error diagnosis}''. However, it is a \textbf{quantum measurement} and therefore \textbf{causes a collapse}. But it \textbf{collapses only the error-related degrees of freedom} (the syndrome subspace), not the logical quantum information encoded in the qubit. Therefore the logical state survives while the error information becomes known.

    \highspace
    For instance, if we have the logical state:
    \begin{equation*}
        \alpha\ket{0} + \beta\ket{1}
    \end{equation*}
    The syndrome does \textbf{not} tell us $\alpha$ or $\beta$, the syndrome only tells us things like ``\emph{error detected: yes}; \emph{location: qubit 5}; \emph{type: bit-flip}''.
\end{definitionbox}

\begin{flushleft}
    \textcolor{Green3}{\faIcon{question-circle} \textbf{How can we perform syndrome measurements without collapsing the logical state?}}
\end{flushleft}
The key insight is that we can perform \textbf{indirect measurements} by coupling the data qubits to an additional qubit, called an \definition{Ancilla Qubit}, and then measuring the ancilla. This way, we can extract information about the error without directly measuring the data qubits, thus preserving the logical quantum state.

\highspace
Usually the ancilla qubit starts in a known state, such as $\ket{0}$. We then \hl{apply a series of quantum gates that entangle the ancilla with the data qubits} in such a way that the state of the ancilla after the interaction contains information about the syndrome.

\begin{examplebox}[: Syndrome Measurement with an Ancilla Qubit]
    We want to know if $q_0$ and $q_1$ are the same or different. We initialize the ancilla qubit in the state $\ket{0}$ and apply a CNOT gate with $q_1$ as the control and the ancilla as the target:
    \begin{center}
        \begin{quantikz}
            \lstick{$q_1$} & \ctrl{1} & \\
            \lstick{Ancilla} & \targ{} &
        \end{quantikz}
    \end{center}
    If $q_1$ is in the state $\ket{0}$, the ancilla remains in $\ket{0}$. If $q_1$ is in the state $\ket{1}$, the ancilla flips to $\ket{1}$. Next, we apply another CNOT gate with $q_0$ as the control and the ancilla as the target:
    \begin{center}
        \begin{quantikz}
            \lstick{$q_0$} & \ctrl{1} & \\
            \lstick{Ancilla} & \targ{} &
        \end{quantikz}
    \end{center}
    After these two gates, the state of the ancilla will be $\ket{0}$ if $q_0$ and $q_1$ are the same (both $\ket{0}$ or both $\ket{1}$) and $\ket{1}$ if they are different. The final value becomes $\text{Ancilla} = q_0 \oplus q_1$, which is the syndrome for this particular check.
\end{examplebox}

\noindent
Finally, we \textbf{measure the ancilla qubit}. The measurement outcome gives the syndrome, which can be used to determine whether an error occurred and, in many cases, which qubit was affected.

\highspace
An important observation is that the syndrome does not exist beforehand. The \hl{ancilla qubit is the physical system that temporarily stores the syndrome information so that it can be measured without directly measuring the data qubits.}

\highspace
\begin{definitionbox}[: Ancilla Qubit]
    An \definition{Ancilla Qubit} is an \textbf{auxiliary qubit} introduced to assist a quantum computation. In QECC, ancilla qubits are used to interact with data qubits, extract information about errors without collapsing the data qubits' states, and help correct errors. They \textbf{do not store the logical quantum information} being protected.
\end{definitionbox}

\highspace
In summary, QECC makes reliable quantum computation possible, but at a \textbf{significant cost}. A logical qubit is not implemented by a single physical qubit. Instead, it requires many physical qubits, including both data qubits and ancilla qubits used for syndrome extraction and error correction.

\highspace
As a consequence, the number of physical qubits required by a quantum computer grows dramatically. In practice, a \textbf{single logical qubit may require tens or even hundreds of physical qubits}. This phenomenon is known as the \definition{Blow-Up Problem}.

\highspace
\begin{definitionbox}[: Quantum Error Correction Code (QECC)]
    A \definition{Quantum Error Correction Code (QECC)} is a \textbf{technique used to protect quantum information} from noise and decoherence. A logical qubit is encoded into multiple physical qubits, introducing redundancy that allows errors to be detected and corrected.

    \highspace
    Error detection is performed through \textbf{syndrome measurements}, typically using \textbf{ancilla qubits}, so that information about errors can be extracted \emph{without} directly measuring the logical quantum state.

    \highspace
    By continuously detecting and correcting errors, QECC enables reliable quantum computations despite the presence of imperfect hardware.
\end{definitionbox}